\section{Fundamentação teórica}
\subsection{Plataformas bilaterais e densidade local}
Rochet e Tirole (2003) demonstram que plataformas bilaterais precisam considerar os efeitos cruzados entre grupos e a distribuição de preços entre os lados. Eisenmann, Parker e Van Alstyne (2006) discutem a necessidade de escolher qual lado subsidiar e como administrar efeitos de rede. A PetCuida conecta tutores e prestadores veterinários; por isso, uma base ampla de tutores sem oferta local suficiente, ou clínicas sem demanda, produz baixo valor para ambos.

Veisdal (2020), em estudo de múltiplos casos de startups de plataformas, examina o problema de iniciar um mercado em que cada lado aguarda a presença do outro. A implicação para a PetCuida é começar com um recorte territorial e uma oferta mínima confiável, em vez de tentar cobrir muitas regiões simultaneamente. A estratégia de densidade local é, portanto, coerente com a teoria e com a limitação de recursos da equipe.

\subsection{Crescimento, escalabilidade e regras de decisão}
Estudo longitudinal sobre uma plataforma bilateral distingue crescimento de escalabilidade e mostra como regras decisórias e efeitos de rede ajudam a transição entre as etapas \cite{platformscaling2023}. Sanasi, Ghezzi e Cavallo (2023) indicam que a escalabilidade em startups tecnológicas continua exigindo experimentação, priorização de hipóteses e exploração das capacidades tecnológicas existentes. A consequência prática é instituir gatilhos de expansão: uma praça só deve ser replicada quando retenção, CAC, margem e oferta estiverem demonstrados.

\subsection{Experimentação e validação}
Felin et al. (2024) discutem a aplicação do método científico às startups e a importância de explicitar teorias e hipóteses testáveis. Leatherbee e Katila (2020) examinam equipes que aplicam práticas Lean Startup e identificam relações entre formulação de hipóteses, sondagem empírica e convergência da ideia de negócio. Werle e Giones (2026) sintetizam a experimentação empreendedora como combinação de descoberta e construção. Para a PetCuida, isso significa não confundir opinião favorável em demonstração guiada com intenção de compra ou comportamento recorrente: cada hipótese deve ter teste, métrica, limiar e decisão. Bocken e Coffay (2023) mostram que a experimentação Lean pode ser combinada com a effectuation, mobilizando recursos e relacionamentos já disponíveis; essa combinação é pertinente ao piloto local com clínicas e organizações comunitárias \cite{bocken2023}.

\subsection{Economia unitária, CAC e retenção}
O custo de aquisição deve ser comparado à contribuição gerada pelo cliente ao longo do relacionamento, considerando também o custo de servir. Estudo brasileiro sobre CLV em empresa SaaS demonstra que canais de aquisição e perfis de clientes afetam o valor de vida do cliente e que aquisição mais cara pode reduzir a rentabilidade \cite{kock2024}. Na PetCuida, o teto inicial de CAC de R\$ 3,29 (seis meses de contribuição de R\$ 0,55) é apenas regra provisória de prudência, não benchmark de mercado. Deve ser recalculado com coortes, retenção observada e margem de contribuição real.